%%%%%%%%%%%%%%%%%%%%%%%%%%%%%%%%%%%%%%%%%%%%%%%%%%%%%%%%%%%%%%%%%%%%%%%%%%%%%%
% Unicode symbols to be used with pdflatex and Sphinx
%%%%%%%%%%%%%%%%%%%%%%%%%%%%%%%%%%%%%%%%%%%%%%%%%%%%%%%%%%%%%%%%%%%%%%%%%%%%%%

% In Emacs `C-u C-x =` gives the symbol information. We use of
% hexadecimal code point.

% Requires `\usepackage[utf8]{inputenc}` with some versions of
% `pdflatex`.

% We use the `textcomp` package because the `sphinx.sty` also uses it.

% You can find the LaTeX command for a Unicode code point in
% http://unicode.scarfboy.com/ or
% https://www.johndcook.com/unicode_latex.html .

% Making upper and lower case script characters in math mode. From
% http://www.peteryu.ca/tutorials/publishing/latex_math_script_styles
% and
% https://tex.stackexchange.com/questions/58098/what-are-all-the-font-styles-i-can-use-in-math-mode.
\DeclareMathAlphabet{\mathpzc}{OT1}{pzc}{m}{it}

%%%%%%%%%%%%%%%%%%%%%%%%%%%%%%%%%%%%%%%%%%%%%%%%%%%%%%%%%%%%%%%%%%%%%%%%%%%%%%
% Please keep the below list ordered by code points!

\DeclareUnicodeCharacter{00AC}{\ensuremath{\neg}}                     % Symbol ¬
\DeclareUnicodeCharacter{00B1}{\ensuremath{\pm}}                      % Symbol ±
\DeclareUnicodeCharacter{00B2}{\ensuremath{^2}}                       % Symbol ²
\DeclareUnicodeCharacter{00B3}{\ensuremath{^3}}                       % Symbol ³
\DeclareUnicodeCharacter{00B7}{\ensuremath{\cdot}}                    % Symbol ·
\DeclareUnicodeCharacter{00B9}{\ensuremath{^1}}                       % Symbol ¹
\DeclareUnicodeCharacter{00D7}{\ensuremath{\times}}                   % Symbol ×
\DeclareUnicodeCharacter{02B0}{\ensuremath{^h}}                       % Symbol ʰ
\DeclareUnicodeCharacter{02B3}{\ensuremath{^r}}                       % Symbol ʳ
\DeclareUnicodeCharacter{02E2}{\ensuremath{^s}}                       % Symbol ˢ
\DeclareUnicodeCharacter{20F0}{\ensuremath{^*}}                       % Symbol  ⃰ 
\DeclareUnicodeCharacter{0302}{\ensuremath{\hat{\phantom{x}}}}        % Symbol ̂
\DeclareUnicodeCharacter{0332}{\ensuremath{\underline{\phantom{x}}}}  % Symbol ̲
\DeclareUnicodeCharacter{0308}{\ensuremath{\ddot{\phantom{x}}}}       % Symbol ̈
\DeclareUnicodeCharacter{0393}{\ensuremath{\Gamma}}                   % Symbol Γ
\DeclareUnicodeCharacter{0394}{\ensuremath{\Delta}}                   % Symbol Δ
\DeclareUnicodeCharacter{0398}{\ensuremath{\Theta}}                   % Symbol Θ
\DeclareUnicodeCharacter{039B}{\ensuremath{\Lambda}}                  % Symbol Λ
\DeclareUnicodeCharacter{039E}{\ensuremath{\Xi}}                      % Symbol Ξ
\DeclareUnicodeCharacter{03A0}{\ensuremath{\Pi}}                      % Symbol Π
\DeclareUnicodeCharacter{03A3}{\ensuremath{\Sigma}}                   % Symbol Σ

% There is not `\Tau` command in LaTeX, so we use `\mathrm{T}`. See
% http://tex.stackexchange.com/questions/1368/why-can-i-only-use-some-capital-greek-letters-inside-my-equations.
\DeclareUnicodeCharacter{03A4}{\ensuremath{\mathrm{T}}}  % Symbol Τ

\DeclareUnicodeCharacter{03A5}{\ensuremath{\Upsilon}}   % Symbol Υ
\DeclareUnicodeCharacter{03A6}{\ensuremath{\Phi}}       % Symbol Φ
\DeclareUnicodeCharacter{03A8}{\ensuremath{\Psi}}       % Symbol Ψ
\DeclareUnicodeCharacter{03A9}{\ensuremath{\Omega}}     % Symbol Ω
\DeclareUnicodeCharacter{03B1}{\ensuremath{\alpha}}     % Symbol α
\DeclareUnicodeCharacter{03B2}{\ensuremath{\beta}}      % Symbol β
\DeclareUnicodeCharacter{03B3}{\ensuremath{\gamma}}     % Symbol γ
\DeclareUnicodeCharacter{03B4}{\ensuremath{\delta}}     % Symbol δ
\DeclareUnicodeCharacter{03B5}{\ensuremath{\epsilon}}   % Symbol ε
\DeclareUnicodeCharacter{03B6}{\ensuremath{\zeta}}      % Symbol ζ
\DeclareUnicodeCharacter{03B7}{\ensuremath{\eta}}       % Symbol η
\DeclareUnicodeCharacter{03B8}{\ensuremath{\theta}}     % Symbol θ
\DeclareUnicodeCharacter{03B9}{\ensuremath{\iota}}      % Symbol ι
\DeclareUnicodeCharacter{03BA}{\ensuremath{\kappa}}     % Symbol κ
\DeclareUnicodeCharacter{03BB}{\ensuremath{\lambda}}    % Symbol λ
\DeclareUnicodeCharacter{03BC}{\ensuremath{\mu}}        % Symbol μ
\DeclareUnicodeCharacter{03BD}{\ensuremath{\nu}}        % Symbol ν
\DeclareUnicodeCharacter{03BE}{\ensuremath{\xi}}        % Symbol ξ
\DeclareUnicodeCharacter{03C0}{\ensuremath{\pi}}        % Symbol π
\DeclareUnicodeCharacter{03C1}{\ensuremath{\rho}}       % Symbol ρ
\DeclareUnicodeCharacter{03C2}{\ensuremath{\varsigma}}  % Symbol ς
\DeclareUnicodeCharacter{03C3}{\ensuremath{\sigma}}     % Symbol σ
\DeclareUnicodeCharacter{03C4}{\ensuremath{\tau}}       % Symbol τ
\DeclareUnicodeCharacter{03C5}{\ensuremath{\upsilon}}   % Symbol υ
\DeclareUnicodeCharacter{03C6}{\ensuremath{\varphi}}    % Symbol φ
\DeclareUnicodeCharacter{03C7}{\ensuremath{\chi}}       % Symbol χ
\DeclareUnicodeCharacter{03C8}{\ensuremath{\psi}}       % Symbol ψ
\DeclareUnicodeCharacter{03C9}{\ensuremath{\omega}}     % Symbol ω
\DeclareUnicodeCharacter{03D1}{\ensuremath{\vartheta}}  % Symbol ϑ

\DeclareUnicodeCharacter{0432}{\ensuremath{\mathrm{PDF\;TODO}}}  % Symbol в
\DeclareUnicodeCharacter{0435}{\ensuremath{\mathrm{PDF\;TODO}}}  % Symbol е
\DeclareUnicodeCharacter{0438}{\ensuremath{\mathrm{PDF\;TODO}}}  % Symbol и
\DeclareUnicodeCharacter{043C}{\ensuremath{\mathrm{PDF\;TODO}}}  % Symbol м
\DeclareUnicodeCharacter{0440}{\ensuremath{\mathrm{PDF\;TODO}}}  % Symbol р
\DeclareUnicodeCharacter{0442}{\ensuremath{\mathrm{PDF\;TODO}}}  % Symbol т
\DeclareUnicodeCharacter{041F}{\ensuremath{\mathrm{PDF\;TODO}}}  % Symbol П
\DeclareUnicodeCharacter{0BA8}{\ensuremath{\mathrm{PDF\;TODO}}}  % Symbol ந
\DeclareUnicodeCharacter{0BBF}{\ensuremath{\mathrm{PDF\;TODO}}}  % Symbol  ி
\DeclareUnicodeCharacter{1100}{\ensuremath{\mathrm{PDF\;TODO}}}  % Symbol ᄀ
\DeclareUnicodeCharacter{11F9}{\ensuremath{\mathrm{PDF\;TODO}}}  % Symbol ᇹ

\DeclareUnicodeCharacter{1D48}{\ensuremath{^d}}  % Symbol ᵈ
\DeclareUnicodeCharacter{1D50}{\ensuremath{^m}}  % Symbol ᵐ
\DeclareUnicodeCharacter{1D56}{\ensuremath{^p}}  % Symbol ᵖ
\DeclareUnicodeCharacter{1D62}{\ensuremath{_i}}  % Symbol ᵢ
\DeclareUnicodeCharacter{1D63}{\ensuremath{_r}}  % Symbol ᵣ
\DeclareUnicodeCharacter{1D64}{\ensuremath{_u}}  % Symbol ᵤ
\DeclareUnicodeCharacter{1D9C}{\ensuremath{^c}}  % Symbol ᶜ


% The command `\text` requires `\usepackage{amsmath}` and the command
% `\textasciiacute` requires `\usepackage{textcomp}`.

\DeclareUnicodeCharacter{2032}{\ensuremath{\text{\textasciiacute}}}  % Symbol ′

\DeclareUnicodeCharacter{2070}{\ensuremath{^0}}  % Symbol ⁰
\DeclareUnicodeCharacter{2074}{\ensuremath{^4}}  % Symbol ⁴
\DeclareUnicodeCharacter{2075}{\ensuremath{^5}}  % Symbol ⁵
\DeclareUnicodeCharacter{2076}{\ensuremath{^6}}  % Symbol ⁶
\DeclareUnicodeCharacter{2077}{\ensuremath{^7}}  % Symbol ⁷
\DeclareUnicodeCharacter{2078}{\ensuremath{^8}}  % Symbol ⁸
\DeclareUnicodeCharacter{2079}{\ensuremath{^9}}  % Symbol ⁹
\DeclareUnicodeCharacter{207A}{\ensuremath{^+}}  % Symbol ⁺
\DeclareUnicodeCharacter{207B}{\ensuremath{^-}}  % Symbol ⁻
\DeclareUnicodeCharacter{207F}{\ensuremath{^n}}  % Symbol ⁿ
\DeclareUnicodeCharacter{2080}{\ensuremath{_0}}  % Symbol ₀
\DeclareUnicodeCharacter{2081}{\ensuremath{_1}}  % Symbol ₁
\DeclareUnicodeCharacter{2082}{\ensuremath{_2}}  % Symbol ₂
\DeclareUnicodeCharacter{2083}{\ensuremath{_3}}  % Symbol ₃
\DeclareUnicodeCharacter{2084}{\ensuremath{_4}}  % Symbol ₄
\DeclareUnicodeCharacter{2085}{\ensuremath{_5}}  % Symbol ₅
\DeclareUnicodeCharacter{2086}{\ensuremath{_6}}  % Symbol ₆
\DeclareUnicodeCharacter{2087}{\ensuremath{_7}}  % Symbol ₇
\DeclareUnicodeCharacter{2088}{\ensuremath{_8}}  % Symbol ₈
\DeclareUnicodeCharacter{2089}{\ensuremath{_9}}  % Symbol ₉
\DeclareUnicodeCharacter{208A}{\ensuremath{_+}}  % Symbol ₊
\DeclareUnicodeCharacter{208B}{\ensuremath{_-}}  % Symbol ₋
\DeclareUnicodeCharacter{2090}{\ensuremath{_a}}  % Symbol ₐ
\DeclareUnicodeCharacter{2091}{\ensuremath{_e}}  % Symbol ₑ
\DeclareUnicodeCharacter{2095}{\ensuremath{_h}}  % Symbol ₕ
\DeclareUnicodeCharacter{2096}{\ensuremath{_k}}  % Symbol ₖ
\DeclareUnicodeCharacter{2097}{\ensuremath{_l}}  % Symbol ₗ
\DeclareUnicodeCharacter{2098}{\ensuremath{_m}}  % Symbol ₘ
\DeclareUnicodeCharacter{2099}{\ensuremath{_n}}  % Symbol ₙ
\DeclareUnicodeCharacter{209A}{\ensuremath{_p}}  % Symbol ₚ
\DeclareUnicodeCharacter{2092}{\ensuremath{_o}}  % Symbol ₒ
\DeclareUnicodeCharacter{209B}{\ensuremath{_s}}  % Symbol ₛ
\DeclareUnicodeCharacter{209C}{\ensuremath{_t}}  % Symbol ₜ
\DeclareUnicodeCharacter{2093}{\ensuremath{_x}}  % Symbol ₓ

% \DeclareUnicodeCharacter{2113}{\ensuremath{\mathpzc{l}}} % Symbol ℓ
\DeclareUnicodeCharacter{2113}{\ensuremath{\ell}} % Symbol ℓ

% The command `\mathbbm` requires `\usepackage{bbm}`.
\DeclareUnicodeCharacter{2115}{\ensuremath{\mathbbm{N}}}  % Symbol ℕ
\DeclareUnicodeCharacter{211A}{\ensuremath{\mathbbm{Q}}}  % Symbol ℚ
\DeclareUnicodeCharacter{211D}{\ensuremath{\mathbbm{R}}}  % Symbol ℝ
\DeclareUnicodeCharacter{2124}{\ensuremath{\mathbbm{Z}}}  % Symbol ℤ

\DeclareUnicodeCharacter{2135}{\ensuremath{\aleph}}           % Symbol ℵ
\DeclareUnicodeCharacter{2190}{\ensuremath{\leftarrow}}       % Symbol ←
\DeclareUnicodeCharacter{2192}{\ensuremath{\rightarrow}}      % Symbol →
\DeclareUnicodeCharacter{2194}{\ensuremath{\leftrightarrow}}  % Symbol ↔

% The command `\twoheadrightarrow` requires `\usepackage{amssymb}`.
\DeclareUnicodeCharacter{21A0}{\ensuremath{\twoheadrightarrow}}  % Symbol ↠

\DeclareUnicodeCharacter{21A6}{\ensuremath{\mapsto}}  % Symbol ↦
\DeclareUnicodeCharacter{21A6}{\ensuremath{\mapsto}}  % Symbol ↦

% The command `\restriction` requires `\usepackage{amssymb}`.
\DeclareUnicodeCharacter{21BE}{\ensuremath{\restriction}}  % Symbol ↾

\DeclareUnicodeCharacter{21D0}{\ensuremath{\Leftarrow}}       % Symbol ⇐
\DeclareUnicodeCharacter{21D2}{\ensuremath{\Rightarrow}}      % Symbol ⇒
\DeclareUnicodeCharacter{21D4}{\ensuremath{\Leftrightarrow}}  % Symbol ⇔

\DeclareUnicodeCharacter{2200}{\ensuremath{\forall}}      % Symbol ∀
\DeclareUnicodeCharacter{2203}{\ensuremath{\exists}}      % Symbol ∃
\DeclareUnicodeCharacter{2204}{\ensuremath{\not\exists}}  % Symbol ∃
\DeclareUnicodeCharacter{2205}{\ensuremath{\emptyset}}    % Symbol ∅
\DeclareUnicodeCharacter{2208}{\ensuremath{\in}}          % Symbol ∈
\DeclareUnicodeCharacter{2209}{\ensuremath{\not\in}}      % Symbol ∉
\DeclareUnicodeCharacter{220B}{\ensuremath{\ni}}          % Symbol ∋
\DeclareUnicodeCharacter{220C}{\ensuremath{\not\ni}}      % Symbol ∌

% The command `\blacksquare` requires `\usepackage{amssymb}`.
\DeclareUnicodeCharacter{220E}{{\tiny \ensuremath{\blacksquare}}} % Symbol ∎

\DeclareUnicodeCharacter{220F}{{\tiny \ensuremath{\prod}}}  % Symbol ∏
\DeclareUnicodeCharacter{2211}{\ensuremath{\sum}}           % Symbol ∑
% 2218
\DeclareUnicodeCharacter{25CB}{\ensuremath{\circ}}          % Symbol ∘
% \DeclareUnicodeCharacter{2219}{\ensuremath{\bullet}}        % Symbol ∙
\DeclareUnicodeCharacter{2022}{\ensuremath{\bullet}}        % Symbol •
\DeclareUnicodeCharacter{221E}{\ensuremath{\infty}}         % Symbol ∞
\DeclareUnicodeCharacter{2223}{\ensuremath{\mid}}           % Symbol ∣
\DeclareUnicodeCharacter{2225}{\ensuremath{\parallel}}      % Symbol ∥
\DeclareUnicodeCharacter{2227}{\ensuremath{\wedge}}         % Symbol ∧
\DeclareUnicodeCharacter{2228}{\ensuremath{\vee}}           % Symbol ∨

\DeclareUnicodeCharacter{2229}{\ensuremath{\cap}}           % Symbol ∩

\DeclareUnicodeCharacter{222A}{\ensuremath{\cup}}           % Symbol ∪
\DeclareUnicodeCharacter{222B}{\ensuremath{\int}}           % Symbol ∫
\DeclareUnicodeCharacter{2234}{\ensuremath{\therefore}}     % Symbol ∴

% The command `\dblcolon` requires `\usepackage{mathtools}`.
\DeclareUnicodeCharacter{2237}{\ensuremath{\dblcolon}} % Symbol ∷

\newcommand{\dotminus}{\mathop{\mbox{$-^{\hspace{-.5em}\cdot}\,$}}}
\DeclareUnicodeCharacter{2238}{\ensuremath{\dotminus}}  % Symbol ∸

\DeclareUnicodeCharacter{223C}{\ensuremath{\sim}}       % Symbol ∼
\DeclareUnicodeCharacter{2243}{\ensuremath{\simeq}}     % Symbol ≃
\DeclareUnicodeCharacter{2245}{\ensuremath{\cong}}      % Symbol ≅
\DeclareUnicodeCharacter{2248}{\ensuremath{\approx}}    % Symbol ≈
\DeclareUnicodeCharacter{2250}{\ensuremath{\doteq}}     % Symbol ≐

\usepackage{mathtools}
% The command `\coloneqq` requires `\usepackage{mathtools}`.
% \DeclareUnicodeCharacter{2254}{\ensuremath{\coloneq}}        % Symbol ≔
\DeclareUnicodeCharacter{2254}{\ensuremath{\coloneqq}}  % Symbol ≔

\newcommand{\eqq}{\stackrel{\text{\tiny ?}}{=}}
\DeclareUnicodeCharacter{225F}{\ensuremath{\eqq}}  % Symbol ≟

\DeclareUnicodeCharacter{2260}{\ensuremath{\neq}}         % Symbol ≠
\DeclareUnicodeCharacter{2261}{\ensuremath{\equiv}}       % Symbol ≡
\DeclareUnicodeCharacter{2262}{\ensuremath{\not\equiv}}   % Symbol ≢
\DeclareUnicodeCharacter{2264}{\ensuremath{\le}}          % Symbol ≤
\DeclareUnicodeCharacter{2265}{\ensuremath{\ge}}          % Symbol ≥
\DeclareUnicodeCharacter{226E}{\ensuremath{\nless}}       % Symbol ≮
\DeclareUnicodeCharacter{226F}{\ensuremath{\ngtr}}        % Symbol ≯
\DeclareUnicodeCharacter{2270}{\ensuremath{\nleq}}        % Symbol ≰
\DeclareUnicodeCharacter{2271}{\ensuremath{\ngeq}}        % Symbol ≱
\DeclareUnicodeCharacter{227A}{\ensuremath{\prec}}        % Symbol ≺
\DeclareUnicodeCharacter{2280}{\ensuremath{\nprec}}        % Symbol ⊀
\DeclareUnicodeCharacter{227C}{\ensuremath{\preceq}}      % Symbol ≼
\DeclareUnicodeCharacter{2282}{\ensuremath{\subset}}      % Symbol ⊂
\DeclareUnicodeCharacter{2284}{\ensuremath{\not\subset}} % Symbol ⊄
\DeclareUnicodeCharacter{2283}{\ensuremath{\supset}}      % Symbol ⊃
\DeclareUnicodeCharacter{2286}{\ensuremath{\subseteq}}    % Symbol ⊆
\DeclareUnicodeCharacter{228E}{\ensuremath{\uplus}}       % Symbol ⊎
\DeclareUnicodeCharacter{2291}{\ensuremath{\sqsubseteq}}  % Symbol ⊑
\DeclareUnicodeCharacter{228F}{\ensuremath{\sqsubset}}  % Symbol ⊏
\DeclareUnicodeCharacter{2293}{\ensuremath{\sqcap}}       % Symbol ⊓
\DeclareUnicodeCharacter{2294}{\ensuremath{\sqcup}}       % Symbol ⊔
\DeclareUnicodeCharacter{2295}{\ensuremath{\oplus}}       % Symbol ⊕
\DeclareUnicodeCharacter{22A2}{\ensuremath{\vdash}}       % Symbol ⊢
\DeclareUnicodeCharacter{22A4}{\ensuremath{\top}}         % Symbol ⊤
\DeclareUnicodeCharacter{22A5}{\ensuremath{\bot}}         % Symbol ⊥
\DeclareUnicodeCharacter{22C0}{\ensuremath{\bigwedge}}    % Symbol ⋀
\DeclareUnicodeCharacter{22C2}{\ensuremath{\bigcap}}      % Symbol ⋂
\DeclareUnicodeCharacter{22C3}{\ensuremath{\bigcup}}      % Symbol ⋃
\DeclareUnicodeCharacter{22EE}{\ensuremath{\vdots}}       % Symbol ⋮
\DeclareUnicodeCharacter{22EF}{\ensuremath{\cdots}}       % Symbol ⋯

% The command `\iddots` requires `\usepackage{mathdots}`.
\DeclareUnicodeCharacter{22F0}{\ensuremath{\iddots}} % Symbol ⋰

\DeclareUnicodeCharacter{230A}{\ensuremath{\lfloor}}  % Symbol ⌊
\DeclareUnicodeCharacter{230B}{\ensuremath{\rfloor}}  % Symbol ⌋
\DeclareUnicodeCharacter{2308}{\ensuremath{\lceil}}   % Symbol ⌈
\DeclareUnicodeCharacter{2309}{\ensuremath{\rceil}}   % Symbol ⌉

\DeclareUnicodeCharacter{2329}{\ensuremath{\langle}}  % Symbol 〈  X〈
\DeclareUnicodeCharacter{232A}{\ensuremath{\rangle}}  % Symbol  〉 X 〉

% The command `\Return` requires `\usepackage{keystroke}`.
\DeclareUnicodeCharacter{23CE}{\ensuremath{\Return}}  % Symbol ⏎

% The command `\text` requires `\usepackage{amsmath}` and the command
% `\ding` requires `\usepackage{pifont}`.
\DeclareUnicodeCharacter{2460}{\ensuremath{\text{\ding{172}}}} % Symbol ①
\DeclareUnicodeCharacter{2461}{\ensuremath{\text{\ding{173}}}} % Symbol ②
\DeclareUnicodeCharacter{2462}{\ensuremath{\text{\ding{174}}}} % Symbol ③
\DeclareUnicodeCharacter{2463}{\ensuremath{\text{\ding{175}}}} % Symbol ④
\DeclareUnicodeCharacter{2464}{\ensuremath{\text{\ding{176}}}} % Symbol ⑤
\DeclareUnicodeCharacter{2465}{\ensuremath{\text{\ding{177}}}} % Symbol ⑥
\DeclareUnicodeCharacter{2466}{\ensuremath{\text{\ding{178}}}} % Symbol ⑦
\DeclareUnicodeCharacter{2467}{\ensuremath{\text{\ding{179}}}} % Symbol ⑧
\DeclareUnicodeCharacter{2468}{\ensuremath{\text{\ding{180}}}} % Symbol ⑨

\DeclareUnicodeCharacter{25C5}{\ensuremath{\triangleleft}}  % Symbol ◅

\DeclareUnicodeCharacter{2660}{\ensuremath{\spadesuit}} % Symbol ♠
\DeclareUnicodeCharacter{266D}{\ensuremath{\flat}}      % Symbol ♭
\DeclareUnicodeCharacter{266F}{\ensuremath{\sharp}}     % Symbol ♯

\DeclareUnicodeCharacter{1D4D0}{\ensuremath{\mathcal{A}}} 
\DeclareUnicodeCharacter{1D4D1}{\ensuremath{\mathcal{B}}}  
\DeclareUnicodeCharacter{1D4D2}{\ensuremath{\mathcal{C}}}  
\DeclareUnicodeCharacter{1D4D3}{\ensuremath{\mathcal{D}}}  
\DeclareUnicodeCharacter{1D4D4}{\ensuremath{\mathcal{E}}}  
\DeclareUnicodeCharacter{1D4D5}{\ensuremath{\mathcal{F}}}  
\DeclareUnicodeCharacter{1D4D6}{\ensuremath{\mathcal{G}}}  
\DeclareUnicodeCharacter{1D4D7}{\ensuremath{\mathcal{H}}}  
\DeclareUnicodeCharacter{1D4D8}{\ensuremath{\mathcal{I}}}  
\DeclareUnicodeCharacter{1D4D9}{\ensuremath{\mathcal{J}}}  
\DeclareUnicodeCharacter{1D4DA}{\ensuremath{\mathcal{K}}}  
\DeclareUnicodeCharacter{1D4DB}{\ensuremath{\mathcal{L}}}  
\DeclareUnicodeCharacter{1D4DC}{\ensuremath{\mathcal{M}}}  
\DeclareUnicodeCharacter{1D4DD}{\ensuremath{\mathcal{N}}}  
\DeclareUnicodeCharacter{1D4DE}{\ensuremath{\mathcal{O}}}  
\DeclareUnicodeCharacter{1D4DF}{\ensuremath{\mathcal{P}}}

\DeclareUnicodeCharacter{1D4E0}{\ensuremath{\mathcal{Q}}}  
\DeclareUnicodeCharacter{1D4E1}{\ensuremath{\mathcal{R}}}  
\DeclareUnicodeCharacter{1D4E2}{\ensuremath{\mathcal{S}}}  
\DeclareUnicodeCharacter{1D4E3}{\ensuremath{\mathcal{T}}}  
\DeclareUnicodeCharacter{1D4E4}{\ensuremath{\mathcal{U}}}  
\DeclareUnicodeCharacter{1D4E5}{\ensuremath{\mathcal{V}}}  
\DeclareUnicodeCharacter{1D4E6}{\ensuremath{\mathcal{W}}}  
\DeclareUnicodeCharacter{1D4E7}{\ensuremath{\mathcal{X}}}  
\DeclareUnicodeCharacter{1D4E8}{\ensuremath{\mathcal{Y}}}  
\DeclareUnicodeCharacter{1D4E9}{\ensuremath{\mathcal{Z}}}  


\DeclareUnicodeCharacter{2713}{\ensuremath{\checkmark}}  % Symbol ✓

% The commands `\llbracket` and `\rrbracket` require
% `\usepackage{stmaryrd}`.
\DeclareUnicodeCharacter{27E6}{\ensuremath{\llbracket}}  % Symbol ⟦
\DeclareUnicodeCharacter{27E7}{\ensuremath{\rrbracket}}  % Symbol ⟧

\DeclareUnicodeCharacter{27E8}{\ensuremath{\langle}}          % Symbol ⟨
\DeclareUnicodeCharacter{27E9}{\ensuremath{\rangle}}          % Symbol 〉
\DeclareUnicodeCharacter{27F6}{\ensuremath{\longrightarrow}}  % Symbol % ⟶

\DeclareUnicodeCharacter{2983}{\ensuremath{\mathrm{PDF\;TODO}}}  % Symbol ⦃
\DeclareUnicodeCharacter{2984}{\ensuremath{\mathrm{PDF\;TODO}}}  % Symbol ⦄
\DeclareUnicodeCharacter{2985}{\ensuremath{\mathrm{PDF\;TODO}}}  % Symbol ⦅
\DeclareUnicodeCharacter{2986}{\ensuremath{\mathrm{PDF\;TODO}}}  % Symbol ⦆

% The commands `\llparenthesis` and `\rrparenthesis` require
% `\usepackage{stmaryrd}`.
\DeclareUnicodeCharacter{2987}{\ensuremath{\llparenthesis}}  % Symbol ⦇
\DeclareUnicodeCharacter{2988}{\ensuremath{\rrparenthesis}}  % Symbol ⦈

% \DeclareUnicodeCharacter{29F5}{\ensuremath{\setminus}}   % Symbol ⧵
\DeclareUnicodeCharacter{2216}{\ensuremath{\setminus}}

\DeclareUnicodeCharacter{2A06}{\ensuremath{\bigcup}}  % Symbol ⋃

\DeclareUnicodeCharacter{D7B0}{\ensuremath{\mathrm{PDF\;TODO}}} % Symbol ힰ

% The command `\mathbbm{1}` requires `\usepackage{bbm}`.
\DeclareUnicodeCharacter{1D7D9}{\ensuremath{\mathbbm{1}}}  % Symbol 𝟙

\DeclareUnicodeCharacter{25C2}{\ensuremath{\triangleleft}}  
\DeclareUnicodeCharacter{25B9}{\ensuremath{\triangleright}}
\DeclareUnicodeCharacter{2286}{\ensuremath{\subseteq}}    % Symbol ⊆
\DeclareUnicodeCharacter{2217}{\em}

\DeclareUnicodeCharacter{1D62}{\ensuremath{_i}}  
\DeclareUnicodeCharacter{2C7C}{\ensuremath{_j}}
\DeclareUnicodeCharacter{2096}{\ensuremath{_k}}

\DeclareUnicodeCharacter{276A}{\mbox{$($}}  
\DeclareUnicodeCharacter{276B}{\mbox{$)$}}  

% Normal: ()
% ❨ ❩
% 2768
% 2769
% ❪ ❫
% 276A
% 276B

\DeclareUnicodeCharacter{261E}{\todo}      % Symbol ☞

\DeclareUnicodeCharacter{1D539}{\mathbb{B}}      % Symbol 𝔹
\DeclareUnicodeCharacter{2102}{\mathbb{C}}      % Symbol ℂ
\DeclareUnicodeCharacter{2098}{_m}      % Symbol ₘ
\DeclareUnicodeCharacter{2115}{\mathbb{N}}      % Symbol ℕ

\DeclareUnicodeCharacter{2713}{\ensuremath{\checkmark}}  % Symbol ✓

\DeclareUnicodeCharacter{25CB}{\ensuremath{\circ}}        % Symbol ○
\DeclareUnicodeCharacter{2A1D}{\ensuremath{\bowtie}}        % Symbol ⨝
\DeclareUnicodeCharacter{22C8}{\ensuremath{\bowtie}}        % Symbol ⋈

\DeclareUnicodeCharacter{2287}{\ensuremath{\supseteq}}        % Symbol ⊇

\DeclareUnicodeCharacter{21DD}{\ensuremath{\rightsquigarrow}}        % Symbol ⇝

\DeclareUnicodeCharacter{1D408}{\ensuremath{\mathbb{I}}} % Symbol 𝐈

\DeclareUnicodeCharacter{2217}{\ensuremath{\ast}}      % Symbol ∗

\DeclareUnicodeCharacter{2298}{\ensuremath{\oslash}}      % Symbol ⊘

